%% CE 49X - Introduction to Computational Thinking and Data Science
%% Midterm Exam - Part 1 ANSWER KEY
%% Weeks 1--7
%% Instructor: Dr. Eyuphan Koc
%% Bogazici University - Spring 2026

\documentclass[11pt,a4paper]{article}

%% ==================== PACKAGES ====================
\usepackage[utf8]{inputenc}
\usepackage[T1]{fontenc}
\usepackage{amsmath,amssymb,amsfonts}
\usepackage{graphicx}
\usepackage[margin=0.7in,headheight=14pt]{geometry}
\usepackage{booktabs}
\usepackage{array}
\usepackage{enumerate}
\usepackage{xcolor}
\usepackage{fancyhdr}
\usepackage{listings}
\usepackage{courier}
\usepackage{tcolorbox}

%% Define custom colors
\definecolor{bunavy}{RGB}{0,20,60}
\definecolor{softblue}{RGB}{100,140,180}
\definecolor{codebg}{RGB}{245,245,245}
\definecolor{answergreen}{RGB}{0,100,0}

%% Listings style for Python
\lstset{
  language=Python,
  basicstyle=\ttfamily\small,
  backgroundcolor=\color{codebg},
  frame=single,
  framerule=0.5pt,
  rulecolor=\color{gray!50},
  keywordstyle=\color{blue!70!black},
  stringstyle=\color{red!60!black},
  commentstyle=\color{green!50!black},
  showstringspaces=false,
  breaklines=true,
  numbers=none,
  xleftmargin=4pt,
  xrightmargin=4pt,
  aboveskip=8pt,
  belowskip=8pt
}

%% Answer box
\newtcolorbox{answerbox}{
  colback=green!5!white,
  colframe=answergreen,
  title=Answer,
  fonttitle=\bfseries
}

%% Set up header and footer
\pagestyle{fancy}
\fancyhf{}
\lhead{CE 49X - Spring 2026}
\chead{Midterm Part 1 -- \textcolor{red}{\textbf{ANSWER KEY}}}
\rhead{Dr. E. Koc}
\lfoot{}
\cfoot{\thepage}
\rfoot{}

%% ==================== DOCUMENT BEGINS ====================
\begin{document}

\begin{center}
\Large\textbf{CE 49X: Introduction to Computational Thinking and Data Science}\\[0.1cm]
\large\textbf{Midterm Exam -- Part 1}\\[0.05cm]
\large\textcolor{red}{\textbf{--- ANSWER KEY ---}}\\[0.2cm]
\end{center}

\noindent\rule{\linewidth}{1pt}
\vspace{0.2cm}

%% ==================== SECTION A: MULTIPLE CHOICE ====================
\section*{Section A: Multiple Choice [20 points]}

\begin{enumerate}
\item \textbf{(b) 4} --- \lstinline{range(10)} produces 0--9; values divisible by 3: \{0, 3, 6, 9\} = 4 elements.

\item \textbf{(c)} \lstinline{import numpy as np} --- This creates the alias \lstinline{np} for the numpy module.

\item \textbf{(b)} \lstinline{array([2, 3, 4])} --- Slicing \lstinline{a[1:4]} returns elements at indices 1, 2, 3 (end index is exclusive).

\item \textbf{(b)} \lstinline{df.fillna(method='ffill')} --- Forward fill propagates the last valid value forward. \lstinline{dropna()} removes rows; \lstinline{interpolate()} computes intermediate values; \lstinline{fillna(0)} fills with zero.

\item \textbf{(d) Position} --- Position on a common scale is the most accurate visual encoding channel for quantitative comparison (Cleveland \& McGill hierarchy).

\item \textbf{(b) Data-ink ratio} --- Tufte's principle: maximize the proportion of ink used to display data vs. non-data elements (gridlines, decoration, chart junk).

\item \textbf{(b) 2.0} --- $z = (66 - 50) / 8 = 16 / 8 = 2.0$.

\item \textbf{(b) Reject $H_0$} --- Since $p = 0.03 < \alpha = 0.05$, we reject the null hypothesis. Note: (a) is a common misconception; p-value is NOT the probability that $H_0$ is true.

\item \textbf{(b) It never identifies the dangerous 5\%} --- This is the ``accuracy trap.'' A model that always predicts the majority class achieves high accuracy but has zero recall for the minority class. It is useless for detecting the cases that matter.

\item \textbf{(b) Features are conditionally independent given the class} --- The ``naive'' assumption is that $P(f_1, f_2, \ldots, f_n | L) = \prod_i P(f_i | L)$. This simplifies computation but is rarely true in practice.
\end{enumerate}

\newpage
%% ==================== SECTION B: WRITTEN QUESTIONS ====================
\section*{Section B: Written Questions [30 points]}

\subsection*{Question 1: Code Tracing \& Debugging [6 points]}

\noindent\textbf{(a) (3 points)}

\begin{answerbox}
The list comprehension filters beams $> 5.0$: \lstinline{[7.2, 12.1, 6.3]}.\\
\lstinline{total = 7.2 + 12.1 + 6.3 = 25.6}, \lstinline{count = 3}.

Output:
\begin{lstlisting}[backgroundcolor=\color{white},frame=none]
Found 3 beams totaling 25.6 m
Average: 8.53 m
\end{lstlisting}

\textbf{Grading:} 1 pt for correct list comprehension result, 1 pt for first line, 1 pt for second line with correct formatting.
\end{answerbox}

\vspace{0.3cm}

\noindent\textbf{(b) (3 points)}

\begin{answerbox}
\textbf{Bug:} The function uses the SPT-N value \lstinline{n} as the dictionary key. When duplicate values appear in the input (5 appears twice, 25 appears twice), the second occurrence overwrites the first. The dictionary ends up with only 3 unique keys: \{5: `soft', 25: `medium', 40: `hard'\}.

\textbf{Current output:} \lstinline{3} (instead of the expected 5).

\textbf{Fix:} Use a list instead of a dict, or use enumerate to create unique keys:
\begin{lstlisting}[backgroundcolor=\color{white},frame=none]
results = []               # Use a list instead
for n in spt_n_values:
    if n < 10:
        results.append('soft')
    ...
\end{lstlisting}

\textbf{Grading:} 1 pt for identifying the bug (dict keys overwrite), 1 pt for stating current output is 3, 1 pt for a valid fix.
\end{answerbox}

\vspace{0.3cm}

\subsection*{Question 2: NumPy \& Pandas [6 points]}

\noindent\textbf{(a) (3 points)}

\begin{answerbox}
A \textbf{view} shares memory with the original array --- modifications to the view affect the original. A \textbf{copy} is an independent array with its own memory.

NumPy slicing creates a \textbf{view}. So \lstinline{subset = speeds[1:4]} is a view of \lstinline{speeds[1]}, \lstinline{speeds[2]}, \lstinline{speeds[3]}. Setting \lstinline{subset[0] = 99} modifies \lstinline{speeds[1]}.

\textbf{Answer:} \lstinline{speeds[1]} is \textbf{99}.

\textbf{Grading:} 1 pt for explaining view vs copy, 1 pt for stating slicing creates a view, 1 pt for correct answer (99).
\end{answerbox}

\vspace{0.3cm}

\noindent\textbf{(b) (3 points)}

\begin{answerbox}
\begin{lstlisting}[backgroundcolor=\color{white},frame=none]
df[(df['AVERAGE_SPEED'] < 20) & (df['NUMBER_OF_VEHICLES'] > 100)][['GEOHASH', 'AVERAGE_SPEED']]
\end{lstlisting}

Also acceptable with \lstinline{.loc}:
\begin{lstlisting}[backgroundcolor=\color{white},frame=none]
df.loc[(df['AVERAGE_SPEED'] < 20) & (df['NUMBER_OF_VEHICLES'] > 100), ['GEOHASH', 'AVERAGE_SPEED']]
\end{lstlisting}

\textbf{Grading:} 1 pt for correct boolean conditions with \lstinline{&}, 1 pt for parentheses around each condition, 1 pt for correct column selection. Deduct 1 pt if using \lstinline{and} instead of \lstinline{&} (won't work on Series).
\end{answerbox}

\newpage
\subsection*{Question 3: Visualization [6 points]}

\noindent\textbf{(a) (3 points)}

\begin{answerbox}
\textbf{Problem 1:} A pie chart with 8 categories makes angular comparison nearly impossible. The 3D perspective further distorts proportions (slices closer to the viewer appear larger). \textbf{Fix:} Use a sorted horizontal bar chart --- position/length encoding is far more accurate than angle.

\textbf{Problem 2:} The jet (rainbow) colormap has no perceptual ordering --- it creates false boundaries and is not colorblind-friendly. \textbf{Fix:} Use a sequential colormap like \texttt{YlOrRd} or \texttt{viridis} where darker/brighter corresponds to higher values.

\textbf{Grading:} 1.5 pts per problem identified + fix. Accept any two valid issues from: too many slices for pie, 3D distortion, rainbow colormap. Accept reasonable alternative fixes.
\end{answerbox}

\vspace{0.3cm}

\noindent\textbf{(b) (3 points)}

\begin{answerbox}
\textbf{Anscombe's Quartet} consists of four datasets that have nearly identical summary statistics (same mean, standard deviation, correlation coefficient, and linear regression line) but look completely different when plotted.

\textbf{Lesson:} Summary statistics alone can hide critical patterns in data. One dataset is linear, one is curved, one has a single outlier pulling the regression, and one has a cluster with an outlier. \textbf{Always visualize your data before drawing conclusions or building models.}

\textbf{Grading:} 1 pt for correct definition (4 datasets, same statistics), 1 pt for stating they look different visually, 1 pt for the lesson (must visualize, statistics can mislead).
\end{answerbox}

\vspace{0.3cm}

\subsection*{Question 4: Statistical Analysis [6 points]}

\noindent\textbf{(a) (3 points)}

\begin{answerbox}
$$z = \frac{x - \mu}{\sigma} = \frac{26.1 - 32.5}{3.2} = \frac{-6.4}{3.2} = -2.0$$

$|z| = 2.0$. By strict inequality $|z| > 2$, this is exactly at the boundary and \textbf{not} extreme. By $|z| \geq 2$, it \textbf{is} extreme.

\textbf{Grading:} 2 pts for correct z-score calculation, 1 pt for correct classification with justification. Accept either interpretation if justified.
\end{answerbox}

\vspace{0.3cm}

\noindent\textbf{(b) (3 points)}

\begin{answerbox}
\textbf{Type I error (false positive):} Rejecting $H_0$ when it is true. Here: concluding the concrete \textbf{does not} meet the 30 MPa specification when it actually does. Result: a good concrete batch is rejected, wasting money and causing delays.

\textbf{Type II error (false negative):} Failing to reject $H_0$ when it is false. Here: concluding the concrete \textbf{meets} the specification when it actually does not. Result: substandard concrete is used in construction.

\textbf{Type II is more dangerous for structural safety} because it leads to using weak concrete in structures, risking collapse and endangering lives. Type I only wastes resources.

\textbf{Grading:} 1 pt for each correct definition with engineering interpretation, 1 pt for identifying Type II as more dangerous with justification.
\end{answerbox}

\newpage
\subsection*{Question 5: Machine Learning \& Naive Bayes [6 points]}

\noindent\textbf{(a) (3 points)}

\begin{answerbox}
For the ``Liquefiable'' class: $TP = 20$, $FP = 15$, $FN = 5$.

$$\text{Precision} = \frac{TP}{TP + FP} = \frac{20}{20 + 15} = \frac{20}{35} \approx 0.571$$

$$\text{Recall} = \frac{TP}{TP + FN} = \frac{20}{20 + 5} = \frac{20}{25} = 0.80$$

\textbf{Prioritize recall.} Missing a liquefiable site (false negative) means building on dangerous ground, risking structural collapse during an earthquake. A false positive only triggers unnecessary ground improvement --- costly but not life-threatening.

\textbf{Grading:} 1 pt for correct precision, 1 pt for correct recall, 1 pt for correct prioritization with engineering justification.
\end{answerbox}

\vspace{0.3cm}

\noindent\textbf{(b) (3 points)}

\begin{answerbox}
Using Bayes' theorem:

$$P(\text{Poor} \,|\, \text{8mm}) = \frac{P(\text{8mm} \,|\, \text{Poor}) \cdot P(\text{Poor})}{P(\text{8mm})}$$

First, compute $P(\text{8mm})$ using the Law of Total Probability:
\begin{align*}
P(\text{8mm}) &= P(\text{8mm} \,|\, \text{Good}) \cdot P(\text{Good}) + P(\text{8mm} \,|\, \text{Poor}) \cdot P(\text{Poor}) \\
&= (0.02)(0.70) + (0.15)(0.30) \\
&= 0.014 + 0.045 = 0.059
\end{align*}

Then:
$$P(\text{Poor} \,|\, \text{8mm}) = \frac{(0.15)(0.30)}{0.059} = \frac{0.045}{0.059} \approx \boxed{0.763}$$

Despite the prior belief that only 30\% of beams are in poor condition, observing an 8mm crack raises this to $\approx 76\%$.

\textbf{Grading:} 1 pt for correct $P(\text{8mm})$ via LOTP, 1 pt for correct application of Bayes' formula, 1 pt for correct final answer ($\approx 0.76$).
\end{answerbox}

\vspace{0.5cm}
\begin{center}
\noindent\rule{0.5\linewidth}{0.4pt}\\[0.3cm]
\textbf{End of Answer Key}
\end{center}

\end{document}
